\section{Conclusions}\label{sec:conclusions}
Side-specific interval and directional disjoint-group penalties
admit explicit original-row proximal calculations in the stated
positive diagonal-metric setting. The box-contraction argument,
scalar threshold and joint dual budget identify the mathematical
conditions needed to recover the intended original objective.
Equivalent lifting, directional row splitting and compact
epigraphs delimit the claim: these objectives were already
expressible, and no priority assertion follows from retaining
the original principal rows.

The completed computational predictions are negative. The
original 456-instance study has a primary 168-observation median
reduction of \OriginalEstimate\%, with 95\% limits \OriginalLower\%
and \OriginalUpper\%. The separate grouped 210-instance study
has a 144-member primary score of \PrimaryEstimate\%, limits
\PrimaryLower\% and \PrimaryUpper\%, and \PrimaryFailed{} retained
failed pairs. Neither establishes the registered 10\% gain.

Correct local mathematics and original-space agreement therefore
do not by themselves establish a faster complete numerical
operation. Future computational work should target a specific
new mechanism and evaluate it prospectively against the strongest
compatible references. The present findings support restricted
derivations, reproducible implementation analysis and complete
negative evidence; they do not establish operational safety,
field benefit or general solver superiority.
